\section{System Implementation}
This section details the concrete implementation of the proposed lightweight Digital Rights Management (DRM) prototype. The system translates the abstract architectural design discussed in Section III into functional production modules, emphasizing low computational footprints and a strict separation of concerns.

\subsection{Software Stack and Module Architecture}
The core ecosystem is engineered entirely within the Python 3.13 runtime environment, leveraging a streamlined stack of asynchronous backend utilities and native multi-threaded libraries. The source configuration is systematically partitioned into three main physical files to maintain decoupled synchronization:
\begin{enumerate}
    \item \textbf{Centralized License Server (`server.py`):} Built upon the asynchronous `FastAPI` framework, this module handles asset cryptographic onboarding, JWT authorization token generation, and the secure retrieval of Content Encryption Keys (CEK) from an in-memory key repository acting as a simulated Hardware Security Module (HSM).
    \item \textbf{Multimedia Packager Module (`packager.py`):} A high-throughput structural handler that manages offline byte operations. It extracts native H.264 transport layers, identifies target framing syntax via Annex-B start-code delimiters, and wraps data arrays inside the `cryptography.hazmat` execution layer.
    \item \textbf{Decryption Client Platform (`client.py`):} An optimized engine that mimics user playback terminal constraints. It communicates via authenticated HTTP headers, extracts the unencrypted DRM metadata block (24 bytes containing the `Content\_ID` and base `IV`), and pipes the raw bits through a selective processing loop.
\end{enumerate}

\subsection{Cryptographic Dispatches and Stream Extraction}
The optimization of the lightweight DRM relies on structural inspection rather than bulk payload mutation. The packager processes the raw H.264 byte stream sequentially, isolating individual Network Abstraction Layer (NAL) units by scanning for standard 3-byte or 4-byte synchronization prefixes (\texttt{0x000001} or \texttt{0x00000001}). Once a boundary is isolated, a logical bitwise masking operation is evaluated against the first byte of the NAL payload:
\begin{equation}
\text{NAL Unit Type} = \text{Header Byte} \land \text{0x1F}
\end{equation}

If the type matches critical configurations, specifically Sequence Parameter Sets (SPS, type 7), Picture Parameter Sets (PPS, type 8), or Instantaneous Decoder Refresh (IDR, type 5), it triggers a localized cryptographic block cycle. The block cipher invokes the Advanced Encryption Standard (AES) operating in Counter (CTR) mode. Nonces are dynamically generated by binding the asset's 8-byte base IV to an 8-byte big-endian representation of the active \texttt{nal\_unit\_type} to guarantee stream alignment. 

The precise implementation of this codec-aware selective encryption loop within the packager component is structured as follows:

\begin{algorithm}[H]
\caption{Codec-Aware Selective Stream Encryption Loop}\label{alg:selective_loop}
\begin{algorithmic}[1]
\For{each $unit$ in $\text{parsed\_nal\_units}$}
    \State $nal\_type \leftarrow unit[\text{"type"}]$
    \State $payload \leftarrow unit[\text{"payload"}]$
    \State $prefix \leftarrow unit[\text{"prefix"}]$
    \State $\text{protected\_stream}.\text{extend}(prefix)$ \Comment{Keep start codes intact}
    \If{$nal\_type = 5$ \textbf{or} $nal\_type = 7$ \textbf{or} $nal\_type = 8$}
        \State $ctr\_nonce \leftarrow base\_iv + \text{to\_bytes}(nal\_type)$
        \State $cipher \leftarrow \text{InitAES\_CTR}(cek, ctr\_nonce)$
        \State $enc\_payload \leftarrow cipher.\text{encrypt}(payload)$
        \State $\text{protected\_stream}.\text{extend}(enc\_payload)$
    \Else
        \State $\text{protected\_stream}.\text{extend}(payload)$ \Comment{Bypass predictive slices}
    \EndIf
\EndFor
\end{algorithmic}
\end{algorithm}

On the client engine end, the reverse operation is handled within \texttt{DRMClientPlayer} via a symmetric real-time pipeline. By evaluating the structural structure array generated during the transport scan, the player invokes low-level OpenSSL subroutines to decipher the encrypted anchoring payloads on the fly while passing standard VCL slices directly to the underlying media buffer without execution penalty.